% !PRINTED 273
\weberpage{273}
Setzt man also
\[
\begin{aligned}
A &= w^2 + w_1^2 + w_2^2,\\
B &= w^2w_1^2 + w^2w_2^2 + w_1^2w_2^2,\\
C &= w\,w_1\,w_2,
\end{aligned}
\]
so erhält man
\[
\begin{array}{cccc}
\omega, & A, & B, & C, \\
-\dfrac{1}{\omega}, & A, & B, & C, \\
\omega+1, & e^{-\frac{\pi i}{12}(n-5)}A, & e^{-\frac{\pi i}{6}(n-5)}B, & C,
\end{array}
\]
so daß zwei dieser Funktionen ebenso wie oben $A$ und $B$ benutzt
werden können. Für $n=5$ zeigt sich aber, daß in den Entwickelungen
von $w$, $w_1$, $w_2$ nach Potenzen von $q$ keine negativen
Potenzen vorkommen und daß also $A$, $B$, $C$ und mithin auch
$w$, $w_1$, $w_2$ selbst konstant sind.

Es läßt sich also die Modulargleichung für $n=5$ in jeder
der drei Formen aufstellen:
\weq{18}{
\begin{aligned}
u_1^2v_2^2 + u_2^2v_1^2 &= 2uv,\\
u_2^2v^2 - v_2^2u^2 &= 2u_1v_1,\\
u_1^2v^2 - v_1^2u^2 &= 2u_2v_2.
\end{aligned}
}

Diese Gleichungen lassen sich auch aus der von \gesperrt{Jacobi}
(Fund. art. 30, gesammelte Werke, Bd. 1, S. 123) gegebenen herleiten.

Wir schließen diese Betrachtungen, indem wir in den einfachsten
Fällen die \gesperrt{Jacobischen Modulargleichungen} aus diesen irrationalen
Formen ableiten.

Wir setzen [vgl. \S 54, (3)]:
\[
\begin{aligned}
\frac{u_2}{u} = x = \sqrt[4]{\varkappa}, &\qquad \frac{u_1}{u} = x' = \sqrt[4]{\varkappa'},\\[6pt]
\frac{v_2}{v} = y = \sqrt[4]{\lambda}, &\qquad \frac{v_1}{v} = y' = \sqrt[4]{\lambda'},
\end{aligned}
\]
% Source scan confirms the radicands \varkappa and \varkappa'.
und eliminieren mittels der Relationen
\weq{19}{
\begin{aligned}
x^8 + x'^8 &= 1, & xx' &= \frac{\sqrt{2}}{u^3}\\
y^8 + y'^8 &= 1, & yy' &= \frac{\sqrt{2}}{v^3}
\end{aligned}
}
% Source scan confirms the denominator v^3, parallel to u^3 above.
die Größen $u$ und $v$.
